\chapter{Kombinatorik og sandsynlighedsteori}
\label{chap:ch4}
Forestil dig, at du skal danne hold på 3 til et semesterprojekt i en klasse med 45 studerende. I første omgang har rækkefølgen, du vælger holdmedlemmerne i, ingen betydning, så du skal kun tænke på kombinationer. Antallet af måder at danne et hold på 3 ud af 45 studerende på er hele 14\,190 muligheder!

Det følgende semester introduceres de studerende til projektstyring med Scrum, hvor hvert hold skal have tre bestemte roller: Scrum Master, Product Owner og Development Team. Denne lille ændring -- at rollerne angives -- forvandler pludselig problemet fra en simpel \textit{kombination} til en \textit{permutation}. Nu springer antallet af mulige måder at fordele rollerne på op til 85\,140!

De 45 studerende er frustrerede over det enorme antal muligheder og holder en fest for at slappe af. Da de er velopdragne, beslutter de, at alle skal give hånd til alle andre præcis én gang. Efter et par minutter regner de det samlede antal håndtryk ud -- 990. De studerende bliver igen overraskede over, hvor hurtigt noget så simpelt som at give hånd kan løbe op.

Som aftenen skrider frem, foreslår en studerende en sjov leg -- en lodtrækning om fem dørpræmier, som alle er forskellige. Med 45 fremmødte studerende og kun fem præmier bliver chancen for ikke at vinde noget bekymrende høje 89 procent. De studerende indser hurtigt, at oddsene ikke er på deres side.

En mindre gruppe er ikke klar til at opgive lykken og beslutter at kaste en mønt 10 gange i håb om at få præcis fem gange krone og vinde spillet. Men da de finder ud af, at sandsynligheden for, at det sker, kun er omkring 25 procent, falder humøret yderligere.

De studerende konkluderer, at i stedet for at stole på tilfældet er det på tide at gå i dybden med kombinatorik og sandsynlighedsteori. Med denne viden kan de bedre forudsige udfald og undgå fremtidige skuffelser både til fester og i projektplanlægningen.

\section{Udfaldsrum og hændelser}
\label{sec:sample_space}
Et \textbf{stokastisk eksperiment }er et eksperiment, der kan føre til forskellige udfald, selv når det gentages under de samme betingelser. Denne tilfældighed er et grundlæggende aspekt af mange ingeniøropgaver.

Tænk på det på denne måde: Lad os sige, at du tester hastigheden af et websted under forskellige forhold. Nogle gange indlæses det hurtigt, andre gange er det langsommere. Selv om du bruger den samme kode og server, får ting som netværkstrafik eller belastning af serveren resultaterne til at variere fra gang til gang. 

\begin{definition}{Stokastisk eksperiment}
    Et stokastisk eksperiment er et eksperiment, der kan give forskellige resultater, selv om du gør alting på samme måde hver gang.
\end{definition}

Eller forestil dig, at du måler signalstyrken i en trådløs enhed. Du får måske lidt forskellige målinger hver gang på grund af ting som interferens eller små ændringer i omgivelserne.

Denne tilfældighed viser sig overalt, fra ydelsestest af software til eksperimenter inden for elektroteknik. Det er vigtigt at forvente den og tage højde for den i sine overvejelser. Ellers risikerer du at træffe beslutninger på grundlag af ufuldstændige eller misvisende data. Når du tager højde for tilfældig variation, kan du lave bedre forudsigelser og design.

For at modellere og analysere et stokastisk eksperiment er det afgørende at kende mængden af mulige udfald. I sandsynlighedsteorien kaldes denne mængde \textbf{udfaldsrummet}, og den betegnes \( S \). Et udfaldsrum kan enten være \textbf{diskret} (bestå af en endelig eller tælleligt uendelig mængde af udfald) eller \textbf{kontinuert} (indeholde et interval af reelle tal). Den præcise definition af et udfaldsrum afhænger ofte af formålet med analysen.

Et \textbf{udfald} er ét muligt resultat af det stokastiske eksperiment, og en \textbf{hændelse} er en vilkårlig delmængde af udfaldsrummet, som kan bestå af ét eller flere udfald. Nedenfor følger nogle eksempler, der illustrerer disse begreber:

\begin{example}[ Netværksforsinkelse] \\
Se på et eksperiment, hvor du måler forsinkelsen (latensen) af datapakker i et netværk. Udfaldsrummet kan defineres ud fra typen af målinger:
\begin{itemize}
    \item Hvis latensen måles som et positivt reelt tal, er udfaldsrummet kontinuert: 
    \[
    S = \{ x \mid x > 0 \}.
    \]
    \item Hvis det vides, at latensen ligger mellem 10 og 100 millisekunder, kan udfaldsrummet indsnævres til:
    \[
    S = \{ x \mid 10 \leq x \leq 100 \}.
    \]
    \item Hvis formålet er at inddele latensen i lav, middel eller høj, bliver udfaldsrummet diskret:
    \[
    S = \{\text{lav}, \text{middel}, \text{høj}\}.
    \]
    \item Ved en simpel vurdering af, om latensen overholder en standardgrænse, kan udfaldsrummet reduceres til:
    \[
    S = \{\text{bestået}, \text{fejlet}\}.
    \]
\end{itemize}
Hvert udfald i disse udfaldsrum repræsenterer én mulig måling af latensen, og hændelser kan defineres som mængder af udfald, fx »latensen er høj«.

\end{example}

Det er grundlæggende i sandsynlighedsregning at forstå, hvad udfaldsrum, udfald og hændelser er, fordi det gør det muligt at definere og arbejde med sandsynligheder for komplekse scenarier i forskellige ingeniørsammenhænge. Lad os opsummere disse centrale begreber:

\begin{definition}{Udfaldsrum, udfald og hændelser}
    \textbf{Udfaldsrum:} Mængden af alle mulige udfald af et stokastisk eksperiment kaldes udfaldsrummet og betegnes \( S \). Udfaldene kan være diskrete eller kontinuerte afhængigt af eksperimentets karakter.
    
    \textbf{Udfald:} Ét muligt resultat af et stokastisk eksperiment.
    
    \textbf{Hændelse:} En vilkårlig delmængde af udfaldsrummet, som kan bestå af ét eller flere udfald.
\end{definition}

\begin{example} Test af en softwareudgivelse \\
    Forestil dig en softwaretestproces, hvor hvert testtilfælde enten kan bestås eller fejle. Udfaldsrummet for et enkelt testtilfælde er diskret og kan skrives som:
    \[
    S = \{\text{bestået}, \text{fejlet}\}.
    \]
    Hvis du kører tre testtilfælde, er det samlede udfaldsrum for alle mulige udfald:
    \begin{align*}
    S = \{&(\text{bestået, bestået, bestået}), (\text{bestået, bestået, fejlet}), (\text{bestået, fejlet, bestået}), \\
    &(\text{bestået, fejlet, fejlet}), (\text{fejlet, bestået, bestået}), (\text{fejlet, bestået, fejlet}), \\
    &(\text{fejlet, fejlet, bestået}), (\text{fejlet, fejlet, fejlet})\}.
    \end{align*}
    Dette udfaldsrum omfatter alle følger af udfald for de tre test.
    
    \begin{itemize}
        \item Det samlede antal mulige udfald er \( 2^3 = 8 \).
        \item En hændelse kunne defineres som »mindst én test fejler«, hvilket ville omfatte udfald som \(\text{(fejlet, bestået, bestået)}\), \(\text{(bestået, fejlet, fejlet)}\) og de øvrige, hvor mindst én test fejler.
    \end{itemize}
    
    \end{example}
    

\begin{example} Komponentkvalitet i produktionen \\
En virksomhed producerer elektroniske komponenter, og hver komponent testes for, om den overholder kvalitetsstandarderne. Testen kan give et af tre udfald: \texttt{bestået}, \texttt{marginal} eller \texttt{fejlet}. Udfaldsrummet er:
\[
S = \{\text{bestået}, \text{marginal}, \text{fejlet}\}.
\]
\textbf{Hændelse:} Antag, at vi er interesseret i hændelsen, at en komponent ikke består kvalitetstesten. Denne hændelse er en mængde af udfald:
\[
E = \{\text{marginal}, \text{fejlet}\}.
\]
\end{example}

Det betyder, at alle de operationer, vi har defineret på mængder, kan overføres direkte til operationer på hændelser:
\begin{itemize}
    \item $A \cup B$: hændelsen, at mindst én af $A$ og $B$ indtræffer.
    \item $A \cap B$: hændelsen, at både $A$ og $B$ indtræffer.
    \item $\skoverline{A}$ eller $A^c$: hændelsen, at $A$ ikke indtræffer.
    \item $A \setminus B$ eller $A - B$: hændelsen, at $A$ indtræffer, men $B$ ikke gør.
    \item $A \Delta B$ eller $A \oplus B$: hændelsen, at enten $A$ eller $B$ indtræffer, men ikke begge (symmetrisk differens).
\end{itemize}

Sandsynlighedsteorien bygger altså direkte på mængdelæren: En sandsynlighed tildeler et numerisk mål til disse delmængder af udfaldsrummet. Senere i kapitlet ser vi, hvordan dette mål defineres og bruges til at ræsonnere om usikkerhed.

Vi kan opsummere overvejelserne i dette afsnit i følgende tabel:

\begin{table}[h]
    \centering
    \renewcommand{\arraystretch}{1.5} % Adjust this value for desired spacing
    \begin{tabular}{|c|c|c|c|}
    \hline
    \textbf{Operation} & \textbf{Boolesk algebra} & \textbf{Logik} & \textbf{Mængdelære} \\
    \hline
    \texttt{NOT} & $\skoverline{x}$ & $\neg x$ & $A^c$ eller $A'$\\
    \hline
    \texttt{OR} & $+$ & $\vee$ & $\cup$ \\
    \hline
    \texttt{AND} & $\cdot$ & $\wedge$ & $\cap$ \\
    \hline
    \texttt{NAND} & $\skoverline{x \cdot y}$ & $\neg (x \wedge y)$ & $(A \cap B)^c$ eller \(\skoverline{A \cap B}\)\\
    \hline
    \texttt{NOR} & $\skoverline{x + y}$ & $\neg (x \vee y)$ & $(A \cup B)^c$ eller \(\skoverline{A \cup B}\) \\
    \hline
    \texttt{XOR} (symmetrisk differens) & $x \oplus y$ & $(x \wedge \neg y) \vee (\neg x \wedge y)$ & $A \triangle B$ eller \(A \oplus B\) \\
    \hline
    Differens & $x \cdot \skoverline{y}$ & $x \wedge \neg y$ & $A - B$ eller \( A \setminus B \) \\
    \hline
    \end{tabular}
    \caption{Sammenligning af operatorer i boolesk algebra, logik og mængdelære}
\end{table}

I denne bog bruger vi notationen $A^c$ for komplementet til mængden $A$, som omfatter alle elementer, der ikke er med i $A$.

\section{Tælleprincipper}

I mange problemer inden for matematik, datalogi og ingeniørfag er det vigtigt at bestemme, på hvor mange måder bestemte hændelser kan indtræffe. Uanset om du ordner elementer, udvælger grupper eller navigerer gennem komplekse scenarier, giver tælleteknikker de grundlæggende værktøjer til at løse disse problemer. Teknikkerne rækker ud over simpel aritmetik og gør det muligt at besvare spørgsmål som:

\begin{itemize}
    \item På hvor mange måder kan vi ordne en mængde af objekter?
    \item Ad hvor mange forskellige veje kan en proces forløbe?
    \item Hvad er sandsynligheden for, at en bestemt hændelse indtræffer, når der er flere muligheder?
\end{itemize}

Tælleteknikker som permutationer, kombinationer og multiplikationsreglen hjælper os med at opgøre disse muligheder systematisk.

\subsection*{Multiplikationsreglen}

Vi begynder med at se på det mest grundlæggende tælleprincip, \textbf{multiplikationsreglen}:

\begin{theorem}[Multiplikationsreglen]
    Lad en handling være beskrevet som en følge af $k$ trin. Antag følgende betingelser:
    
    \begin{itemize}
        \item Der er $n_1$ måder at udføre trin 1 på.
        \item Der er $n_2$ måder at udføre trin 2 på for hver måde at udføre trin 1 på.
        \item Der er $n_3$ måder at udføre trin 3 på for hver mulig udførelse af de to første trin, og tilsvarende for hvert af de efterfølgende trin.
    \end{itemize}
    
    Så er det samlede antal måder at udføre hele handlingen på givet ved:
    
    \[
    n_1 \times n_2 \times \cdots \times n_k.
    \]
\end{theorem}
    
\begin{example}
    Antag, at du skal vælge et måltid på en restaurant. Du har følgende muligheder:
        
    \begin{itemize}
            \item 3 valgmuligheder til hovedretten.
            \item 4 valgmuligheder til tilbehøret.
            \item 2 valgmuligheder til drikkevaren.
    \end{itemize}
        
    Ifølge multiplikationsreglen er det samlede antal måder at vælge et måltid på:
        
    \[
        3 \times 4 \times 2 = 24.
    \]
        
    Der er altså 24 forskellige måltidskombinationer at vælge imellem.
\end{example}

\begin{example} Tilvalg til biler

En bilproducent leverer biler med udvalgte
tilvalg. Hver bil bestilles

\begin{itemize}
    \item med eller uden automatgear
    \item med eller uden soltag
    \item med ét ud af tre valgmuligheder for stereoanlæg
    \item med én ud af fire udvendige farver
\end{itemize}

Hvis udfaldsrummet består af mængden af alle mulige biltyper, hvad er så antallet af udfald i udfaldsrummet?
\end{example}

\begin{solution}
Ved hjælp af multiplikationsreglen kan vi beregne det samlede antal mulige biltyper ved at gange antallet af valgmuligheder for hvert tilvalg:

\begin{itemize}
    \item 2 valgmuligheder for gearkassen (med eller uden automatgear)
    \item 2 valgmuligheder for soltaget (med eller uden soltag)
    \item 3 valgmuligheder for stereoanlægget
    \item 4 valgmuligheder for den udvendige farve
\end{itemize}

Det samlede antal mulige biltyper er derfor:

\[
2 \times 2 \times 3 \times 4 = 48
\]

Der er altså 48 forskellige mulige biltyper i udfaldsrummet.
\end{solution}

\subsection*{Tilbagelægning og rækkefølge ved optælling}

Dernæst ser vi på en vigtig skelnen i tælleprincipperne, nemlig mellem optælling med og uden tilbagelægning:

\begin{definition}[Optælling med og uden tilbagelægning]
    Når vi tæller antallet af måder at udvælge objekter fra en mængde på, er der to almindelige scenarier:

    \begin{itemize}
        \item \textbf{Med tilbagelægning}: Et objekt kan udvælges mere end én gang.
        \item \textbf{Uden tilbagelægning}: Når et objekt er udvalgt, kan det ikke vælges igen.
    \end{itemize}
    \end{definition}
    
\begin{example}
    Antag, at du har en pose med 5 kugler i forskellige farver. Du trækker 2 kugler:
\begin{itemize}
    \item \textbf{Med tilbagelægning}: Den første kugle lægges tilbage i posen, før den anden trækkes. Der er \(5 \times 5 = 25\) mulige udfald.
    \item \textbf{Uden tilbagelægning}: Den første kugle lægges ikke tilbage, så antallet af udfald er \(5 \times 4 = 20\).
\end{itemize}
\end{example}

\begin{example}
    Tre personer trækker kort efter tur fra et almindeligt spil kort med 52 kort. Målet er at finde spar es. Lad os se på de to scenarier: med tilbagelægning og uden tilbagelægning.

    \textbf{Uden tilbagelægning}

    I dette scenarie lægges et trukket kort ikke tilbage i bunken, så det samlede antal tilgængelige kort bliver mindre efter hver trækning.
    \begin{itemize}
        \item Første trækning: Den første person har en chance på $\frac{1}{52}$ for at trække spar es.
        \item Anden trækning: Hvis den første person ikke trækker spar es, er der nu 51 kort tilbage, og den anden person har en chance på $\frac{1}{51}$ for at trække spar es.
        \item Tredje trækning: Hvis spar es ikke er trukket af de to første personer, har den tredje person en chance på $\frac{1}{50}$ for at trække det.
    \end{itemize}
    
    Sandsynlighederne ændrer sig ved hver trækning, fordi det samlede antal kort bliver mindre, og tidligere trukne kort ikke er tilgængelige.
    
    \textbf{Med tilbagelægning}

    I dette scenarie lægges hvert trukket kort tilbage i bunken, som blandes, før den næste person trækker. Det samlede antal kort er dermed konstant.
    \begin{itemize}
        \item Første trækning: Den første person har en chance på $\frac{1}{52}$ for at trække spar es.
        \item Anden trækning: Da kortet lægges tilbage og blandes ind i bunken igen, har den anden person også en chance på $\frac{1}{52}$ for at trække spar es.
        \item Tredje trækning: Tilsvarende har den tredje person en chance på $\frac{1}{52}$ for at trække spar es.
    \end{itemize}
    
    Sandsynlighederne er de samme ved hver trækning, fordi bunken bringes tilbage til sin oprindelige tilstand efter hver trækning.

\end{example}

\begin{example}
    Forestil dig, at du har en gruppe på 5 studerende: Alice, Bob, Charlie, David og Eve. Du skal udvælge 2 af dem i forskellige scenarier, som illustrerer, hvornår rækkefølgen af udvælgelsen har betydning, og hvornår den ikke har.
    \begin{itemize}
        \item \textbf{Rækkefølgen har betydning}: At vælge Alice som anfører og Bob som viceanfører er et andet udfald end at vælge Bob som anfører og Alice som viceanfører.
        \end{itemize}
    Forestil dig nu, at du blot udvælger 2 studerende til en læsegruppe uden bestemte roller. Her har rækkefølgen ingen betydning.
    \begin{itemize}
        \item \textbf{Rækkefølgen har ingen betydning:}  At vælge Alice og Bob betragtes som det samme udfald som at vælge Bob og Alice; der er ingen forskel på de to rækkefølger, da der ikke er tildelt roller.
    \end{itemize}
    \label{ex:order_matters}
\end{example}

\autoref{ex:order_matters} illustrerer forskellen mellem \textit{permutationer} og \textit{kombinationer}, to grundlæggende tælleprincipper, som er meget udbredte i sandsynlighedsteori og kombinatorik.

\begin{definition}[Permutation og kombination]
    \begin{itemize}
        \item \textbf{Permutation (rækkefølgen har betydning)}: Forskellige rækkefølger tælles som forskellige udfald, hvilket giver et større antal.
        \item \textbf{Kombination (rækkefølgen har ingen betydning)}: Rækkefølger behandles som ens, hvilket giver et mindre antal.
    \end{itemize}
\end{definition}

Vi ser først på permutationer, som bruges, når rækkefølgen af udvælgelsen har betydning.

\subsection*{Permutationer}

Se på en mængde af elementer, fx $S=\{a, b, c\}$. En permutation af elementerne er en ordnet følge af elementerne. For eksempel er $a b c, a c b, b a c, b c a, c a b$ og $c b a$ alle permutationerne af elementerne i $S$.

\begin{proposition}[Permutationer af \(n\) forskellige objekter]
    Antallet af ordnede opstillinger (permutationer) af $n$ forskellige objekter er
    \[
    n! = n \cdot (n-1) \cdot \dots \cdot 2 \cdot 1.
    \]
\end{proposition}

Resultatet er en direkte anvendelse af multiplikationsreglen. For at danne en permutation begynder du med at vælge et element til den første plads blandt alle $n$ elementer. Derefter vælger du et element til den anden plads blandt de resterende $n-1$ elementer, så til den tredje plads blandt de resterende $n-2$ elementer og fortsætter på denne måde, indtil alle pladser er besat. Sådanne opstillinger kaldes ofte lineære permutationer.

\begin{example}
    Man siger, at enhver blanding af et spil kort kun er forekommet én gang i historien. Det skyldes, at antallet af måder at blande et spil med 52 kort på er \(52!\), hvilket er et astronomisk stort tal.

    \[52! \approx 8.07 \times 10^{67}\]
\end{example}

\begin{example}
    Antag, at du har 5 forskellige bøger på en hylde. Du vil stille dem i en anden rækkefølge. Antallet af måder at stille bøgerne på er
    \[5! = 5 \times 4 \times 3 \times 2 \times 1 = 120\]
\end{example}

Der er tilfælde, hvor vi kun er interesseret i at ordne en delmængde af elementerne fra en større mængde. Formlen for antallet af sådanne opstillinger følger også af multiplikationsreglen.

\begin{theorem}[Permutationer af delmængder]
    For hele tal $n \ge r \ge 0$ er antallet af ordnede udvælgelser af $r$ forskellige objekter blandt $n$ forskellige objekter
    \[
    P_r^n = P(n,r) = n \cdot (n-1) \cdots (n-r+1) = \frac{n!}{(n-r)!}.
    \]
\end{theorem}

\begin{example}
    Antag, at du har 5 forskellige bøger på en hylde, og at du vil stille 3 af dem i en anden rækkefølge. Antallet af måder at stille de 3 bøger på er
    \[
    P_3^5=5 \times 4 \times 3 = 60
    \]
\end{example}

\begin{example}
    Der er 10 deltagere i en konkurrence. Kun tre vinder, nemlig 1., 2. eller 3. præmie. Hvad er de mulige resultater?
\end{example}
\begin{solution}
    Antallet af måder at uddele præmierne på er antallet af permutationer af 3 objekter udvalgt blandt 10, dvs.
    \[
    P_3^{10}= \frac{10!}{(10-3)!}=\frac{10!}{7!} = 10 \times 9 \times 8 = 720
    \]
    Der er altså 720 mulige udfald af præmieuddelingen.
\end{solution}

\subsection*{Kombinationer}

Når rækkefølgen af udvælgelsen ikke har betydning, bruger vi begrebet kombinationer. Kombinationer bruges, når vi er interesseret i at udvælge en delmængde af elementer fra en større mængde uden hensyn til den rækkefølge, de udvælges i. Lad os begynde med et par eksempler, der illustrerer begrebet kombinationer.

\begin{example}
    Antag, at du har en gruppe på 5 studerende: Alice, Bob, Charlie, David og Eve. Du skal udvælge 2 af dem til at danne en læsegruppe. Rækkefølgen, du vælger de studerende i, har ingen betydning. De mulige kombinationer er:
    \begin{itemize}
        \item Alice og Bob
        \item Alice og Charlie
        \item Alice og David
        \item Alice og Eve
        \item Bob og Charlie
        \item Bob og David
        \item Bob og Eve
        \item Charlie og David
        \item Charlie og Eve
        \item David og Eve
    \end{itemize}
    Rækkefølgen af de studerende i læsegruppen har ingen betydning, så kombinationerne betragtes som ens.
\end{example}

\begin{example}
    Maria har tre billetter til en koncert. Hun vil selv bruge den ene billet. De to andre billetter kan hun så tilbyde til hvilke som helst af fire venner (Ann, Beth, Chris, Dave). På hvor mange måder kan 2 personer udvælges blandt 4 til at gå til koncert?
\end{example}

\begin{example}
Et printkort har fire forskellige pladser, hvor en komponent kan placeres. Hvis tre ens komponenter skal placeres på kortet, hvor mange forskellige designs er så mulige?
\end{example}

\begin{solution}
    Da der kun kan placeres én komponent på hver plads, begrænser placeringen af en komponent på en plads straks valgmulighederne for den næste komponent.
        
    \begin{enumerate}
        \item Besæt plads 1, 2 og 3.
        \item Besæt plads 1, 2 og 4.
        \item Besæt plads 1, 3 og 4.
        \item Besæt plads 2, 3 og 4.
    \end{enumerate}
\end{solution}

Disse eksempler illustrerer begrebet kombinationer, hvor rækkefølgen af udvælgelsen ikke har betydning. Formlen for kombinationer udledes af formlen for permutationer ved at dividere med antallet af måder at ordne de $r$ elementer på.

\begin{theorem}[Kombinationer]
    Antallet af kombinationer af $r$ elementer udvalgt fra en mængde med $n$ forskellige elementer er givet ved
    \[
    C_r^n = \binom{n}{r} = \frac{n!}{r!(n-r)!}
    \]
\end{theorem}

Dette kaldes også \textbf{binomialkoefficienten}, som betegnes $\binom{n}{r}$ og læses »n over r«. Den kaldes binomialkoefficienten, fordi den optræder i binomialsætningen, som udvikler potenser af et binomium (et toleddet udtryk):

\begin{theorem}[Binomialsætningen]
    I algebra bruges binomialkoefficienten til at udvikle potenser af binomier. Ifølge binomialsætningen gælder

$$
(a+b)^n=\sum_{k=0}^n\binom{n}{k} a^k b^{n-k}
$$

\end{theorem}

Sætningen siger, at udviklingen af binomiet $(a+b)^n$ er summen af leddene $\binom{n}{k} a^k b^{n-k}$ for $k=0,1,2,\ldots,n$. Binomialkoefficienten $\binom{n}{k}$ angiver antallet af måder at vælge $k$ elementer fra en mængde med $n$ elementer på. Vi vil ikke lægge yderligere vægt på binomialsætningen her, men den er et grundlæggende begreb i algebra og kombinatorik og bruges i vid udstrækning i sandsynlighedsteori.

Vi afslutter gennemgangen af tælleprincipper med principperne for optælling med tilbagelægning.

\begin{proposition}[Med tilbagelægning]
Ved udvælgelse blandt $n$ typer gælder:
\begin{itemize}
    \item \textbf{Ordnet med tilbagelægning}: $n^r$ udfald.
    \item \textbf{Uordnet med tilbagelægning}: $\displaystyle \binom{n+r-1}{r}$ udfald. % (stars and bars)
\end{itemize}
\end{proposition}

\section{Grundlæggende sandsynlighedsregning}
Sandsynlighed er et mål for, hvor sandsynligt eller hvor stor chancen er for, at hændelser i forbindelse med et stokastisk eksperiment indtræffer. Når du fx hører »Chancen for regn i dag er 30\%«, udtrykker det usikkerhed om hændelsen, at det regner. Sandsynligheder ligger mellem 0 og 1 (eller tilsvarende mellem 0\% og 100\%). En umulig hændelse har sandsynlighed 0, og hele udfaldsrummet har sandsynlighed 1. I generelle sandsynlighedsmodeller kan en hændelse dog have sandsynlighed 0 uden at være tom eller sandsynlighed 1 uden at indeholde alle mulige udfald; sådanne hændelser kaldes henholdsvis næsten umulige og næsten sikre.

Sandsynligheder kan fortolkes på forskellige måder:

\begin{itemize}
    \item \textbf{Objektiv (eller klassisk) sandsynlighed}: Denne tilgang, der ofte kaldes klassisk sandsynlighed, bruges, når udfaldene er lige sandsynlige, fx ved kast med en fair terning eller en mønt. Sandsynlighederne tildeles ud fra antagelsen om, at alle udfald har lige stor chance for at indtræffe. Når man fx kaster en fair sekssidet terning, er sandsynligheden for at slå en 3'er \( \frac{1}{6} \), fordi der er 6 lige sandsynlige udfald (1, 2, 3, 4, 5, 6), og kun ét af dem er en 3'er. Sandsynligheden er den samme for alle iagttagere.
        
    \item \textbf{Relativ hyppighed (empirisk sandsynlighed)}: Empirisk sandsynlighed bygger på observationer fra eksperimenter frem for teoretiske beregninger. Hvis en softwaretester fx kører en belastningstest på en server 100 gange, og serveren går ned 7 gange, er den empiriske sandsynlighed for et nedbrud \( \frac{7}{100} = 0.07 \). Denne tilgang bygger på faktiske data frem for antagelser eller intuition.
        
    \item \textbf{Subjektiv sandsynlighed}: Den afspejler vores personlige overbevisning eller grad af tiltro til et udfald. Forskellige personer kan tildele den samme hændelse forskellige sandsynligheder ud fra deres viden eller synsvinkel. Du og dine venner diskuterer Danmarks chancer for at vinde VM. Ud fra holdets seneste præstationer og styrke vurderer du chancen til 10\%. En mere optimistisk ven sætter den derimod til 20\%, mens en anden kun giver 5\%, fordi der er stærkere konkurrenter. Det illustrerer subjektiv sandsynlighed, hvor hver persons vurdering varierer ud fra personlige overbevisninger og forudindtagethed frem for objektive data.
    

\end{itemize}

Når vi tildeler sandsynligheder, er det afgørende, at summen af alle sandsynligheder i et eksperiment er 1. Det sikrer overensstemmelse med fortolkningen som relativ hyppighed.

Vi begynder med at opstille sandsynlighedsregningens aksiomer, som danner grundlaget for, hvordan sandsynligheder tildeles hændelser. Aksiomerne fastlægger de grundlæggende egenskaber, som ethvert sandsynlighedsmål skal opfylde.
    
\begin{axiom}[Sandsynlighedsregningens aksiomer]
    \begin{itemize}
        \item \textbf{Aksiom 1:} For enhver hændelse \(A\) gælder \( 0 \leq P(A) \leq 1 \).
        \item \textbf{Aksiom 2:} Sandsynligheden for udfaldsrummet \(S\) er \(P(S) = 1\).
        \item \textbf{Aksiom 3:} Hvis \(A_1, A_2, A_3, \cdots\) er disjunkte hændelser, så gælder \(P\left(A_1 \cup A_2 \cup A_3 \cdots\right) = P\left(A_1\right) + P\left(A_2\right) + P\left(A_3\right) + \cdots\)
    \end{itemize}
\end{axiom}

Egenskaben \( 0 \leq P(A) \leq 1 \) svarer til kravet om, at en relativ hyppighed skal ligge mellem 0 og 1. Egenskaben \( P(S) = 1 \) er en følge af, at der ved hvert forsøg i et eksperiment indtræffer et udfald fra udfaldsrummet. Derfor er den relative hyppighed af \( S \) lig med 1. Egenskab 3 indebærer, at hvis hændelserne \( A_1 \) og \( A_2 \) ikke har nogen udfald til fælles, er den relative hyppighed af udfald i \( A_1 \cup A_2 \) summen af de relative hyppigheder af udfaldene i \( A_1 \) og \( A_2 \).

I de næste afsnit ser vi nærmere på sandsynligheden for hændelser, og hvordan den beregnes.

\subsection*{Sandsynligheden for en hændelse}
Sandsynligheden for en hændelse er et mål for, hvor sandsynligt det er, at hændelsen indtræffer. Den betegnes \( P(A) \), hvor \( A \) er hændelsen, og den ligger mellem 0 og 1. Den tomme hændelse har sandsynlighed 0, og udfaldsrummet har sandsynlighed 1.

\begin{definition}[Klassisk sandsynlighed for lige sandsynlige udfald]
    Hvis et endeligt udfaldsrum består af lige sandsynlige udfald, er sandsynligheden for en hændelse \( A \) forholdet mellem antallet af udfald i \(A\) og det samlede antal udfald i udfaldsrummet:
    \[
    P(A) = \frac{\text{Antal gunstige udfald}}{\text{Samlet antal udfald}}
    \]
\end{definition}

\begin{example}
    Antag, at du tester et softwaremodul med 10 forskellige testtilfælde. Det vides, at 3 af testtilfældene fejler på grund af en fejl i koden. Hvis du vælger et tilfældigt testtilfælde og kører det, hvad er så sandsynligheden for, at det valgte testtilfælde fejler?
\end{example}

\begin{solution}
    Her er hændelsen $A$ »testtilfældet fejler«.
\begin{itemize}
    \item Antal gunstige udfald (fejlende testtilfælde) $=3$
    \item Samlet antal udfald (testtilfælde i alt) $=10$
\end{itemize}

Ved hjælp af formlen får vi:

$$
P(A)=\frac{\text { Antal gunstige udfald }}{\text { Samlet antal udfald }}=\frac{3}{10}=0.3
$$


Sandsynligheden for, at et tilfældigt valgt testtilfælde fejler, er altså 0.3, hvilket betyder, at der er $30 \%$ chance for fejl.

\end{solution}

\begin{example}
    Forestil dig et softwareudviklingsmiljø med 50 filer, som består af 20 Python-scripts, 15 Java-filer og 15 konfigurationsfiler. Hvis du vælger en tilfældig fil at redigere, hvad er så sandsynligheden for, at filen er et Python-script?
\end{example}

\begin{solution}
Her er hændelsen $A$ »den valgte fil er et Python-script«.
\begin{itemize}
    \item Antal gunstige udfald (Python-scripts) $=20$
    \item Samlet antal udfald (filer i alt) $=50$
\end{itemize}

Ved hjælp af formlen får vi:

$$
P(A)=\frac{\text { Antal gunstige udfald }}{\text { Samlet antal udfald }}=\frac{20}{50}=0.4
$$

Sandsynligheden for at vælge et Python-script er altså 0.4, hvilket betyder, at der er $40 \%$ chance for at vælge en Python-fil blandt filerne.

\end{solution}

\section{Sandsynlighed for sammensatte hændelser og mængdeoperationer}
Sammensatte hændelser dannes ved at anvende de grundlæggende mængdeoperationer på enkelte hændelser. Vi møder ofte foreningsmængder af hændelser som $A \cup B$, fællesmængder af hændelser som $A \cap B$ og komplementer af hændelser som $A^c$. Disse sammensatte hændelser er ofte af særlig interesse, og deres sandsynligheder kan tit udledes af sandsynlighederne for de enkelte hændelser, de er sammensat af. Det er vigtigt at forstå disse mængdeoperationer for at kunne beregne sandsynligheden for sammensatte hændelser korrekt. I dette afsnit ser vi på, hvordan foreningsmængder af hændelser og andre mængdeoperationer kan bruges til at bestemme sandsynlighederne for mere komplekse hændelser.

Når vi arbejder med hændelser, svarer fællesmængden til \texttt{AND}, mens foreningsmængden svarer til \texttt{OR}. Sandsynligheden for fællesmængden af hændelserne $A$ og $B$, som betegnes $P(A \cap B)$, kan også skrives $P(A, B)$ eller $P(A B)$.

Ud fra sandsynlighedsregningens aksiomer kan vi udlede følgende regler for sandsynligheder:

\begin{theorem}{Regler for sandsynlighed}
    \begin{itemize}
        \item \textbf{Komplementreglen:} Sandsynligheden for komplementet til hændelse $A$ er
    
        \[
        P(A^c) = 1 - P(A)
        \]

        \item \textbf{Reglen for den tomme mængde:} Sandsynligheden for den tomme mængde er 0, dvs.
        
        \[P(\emptyset) = 0\]
        
        \item \textbf{Additionsreglen:} For to vilkårlige hændelser $A$ og $B$ er sandsynligheden for foreningsmængden af hændelserne $A$ og $B$ givet ved
        \[
        P(A \cup B) = P(A) + P(B) - P(A \cap B)
        \]

        \item \textbf{Differensreglen:} Sandsynligheden for differensen mellem hændelserne $A$ og $B$ er givet ved
        \[
        P(A-B)=P(A)-P(A \cap B)
        \]
        \item \textbf{Delmængdereglen:} Hvis $A$ er en delmængde af $B$ ($A \subset B$), så gælder
        \[P(A) \leq P(B)\].
    \end{itemize}
\end{theorem}

Vi kan udlede komplementreglen ved at bemærke følgende:

\begin{align*}
    1 &= P(S) & \text{(aksiom 2)} \\
      &= P(A \cup A^c) & \text{(definitionen af komplement)} \\
      &= P(A) + P(A^c) & \text{(da } A \text{ og } A^c \text{ er disjunkte)}
\end{align*}

Da $\emptyset=S^c$, kan vi bruge komplementreglen til at slutte, at $P(\emptyset)=1-P(S)=0$. Det er intuitivt, fordi en hændelse pr. definition indtræffer, når udfaldet af det stokastiske eksperiment er en del af hændelsen. Men da den tomme mængde ikke indeholder nogen elementer, kan intet udfald af eksperimentet nogensinde tilhøre den, og derfor er dens sandsynlighed nul.

Differensreglen kan udledes ved at vise, at $P(A)=P(A \cap B)+P(A-B)$. Bemærk, at de to mængder $A \cap B$ og $A-B$ er disjunkte, og at deres foreningsmængde er $A$. Ifølge sandsynlighedsregningens tredje aksiom gælder derfor

\begin{align*}
    P(A) &= P\left((A \cap B) \cup (A - B)\right) & \text{(da } A = (A \cap B) \cup (A - B)) \\
         &= P(A \cap B) + P(A - B) & \text{(da } A \cap B \text{ og } A - B \text{ er disjunkte)}
    \end{align*}
    

Additionsreglen får vi ved at bemærke, at $A$ og $B-A$ er disjunkte mængder, og at deres foreningsmængde er $A \cup B$. Dermed gælder

\begin{align*}
    P(A \cup B) &= P(A \cup (B - A)) && (\text{da } A \cup B = A \cup (B - A)) \\
                &= P(A) + P(B - A) && (\text{da } A \text{ og } B - A \text{ er disjunkte}) \\
                &= P(A) + P(B) - P(A \cap B) && (\text{jf. differensreglen})
    \end{align*}
    
Endelig er delmængdereglen en direkte følge af, at hvis $A \subset B$, kan $B$ skrives som foreningsmængden af $A$ og $B-A$. Da $A$ og $B-A$ er disjunkte, har vi $P(B)=P(A)+P(B-A) \geq P(A)$.

Vi afslutter afsnittet med nogle eksempler, der viser, hvordan disse regler anvendes til at beregne sandsynligheder for sammensatte hændelser.

\begin{example}
    En virksomhed har afgivet bud på to store byggeprojekter. Virksomhedens direktør vurderer, at sandsynligheden for at vinde den første kontrakt er 0.6, sandsynligheden for at vinde den anden kontrakt er 0.4, og sandsynligheden for at vinde begge kontrakter er 0.2.

    \begin{enumerate}[label=(\alph*)]
        \item Hvad er sandsynligheden for, at virksomheden vinder mindst én kontrakt?
        \item Hvad er sandsynligheden for, at virksomheden vinder den første kontrakt, men ikke den anden?
        \item Hvad er sandsynligheden for, at virksomheden ikke vinder nogen af kontrakterne?
        \item Hvad er sandsynligheden for, at virksomheden vinder præcis én kontrakt?
    \end{enumerate}
\end{example}

\begin{solution}
    Lad $A$ være hændelsen, at virksomheden vinder den første kontrakt, og $B$ hændelsen, at virksomheden vinder den anden kontrakt. Vi har givet:
    \begin{itemize}
        \item $P(A) = 0.6$
        \item $P(B) = 0.4$
        \item $P(A \cap B) = 0.2$
    \end{itemize}
    
    \begin{enumerate}[label=(\alph*)]
        \item Sandsynligheden for, at virksomheden vinder mindst én kontrakt, er sandsynligheden for foreningsmængden af hændelserne $A$ og $B$. Ifølge additionsreglen er:
        
        \begin{align*}
            P(A \cup B) &= P(A) + P(B) - P(A \cap B) \\
                        &= 0.6 + 0.4 - 0.2 \\
                        &= 0.8
        \end{align*}
        
        Sandsynligheden for, at virksomheden vinder mindst én kontrakt, er altså 0.8.
        
        \item Sandsynligheden for, at virksomheden vinder den første kontrakt, men ikke den anden, er sandsynligheden for differensen mellem hændelserne $A$ og $B$. Ifølge differensreglen er:
        
        \begin{align*}
            P(A - B) &= P(A) - P(A \cap B) \\
                    &= 0.6 - 0.2 \\
                    &= 0.4
        \end{align*}
        
        Sandsynligheden for, at virksomheden vinder den første kontrakt, men ikke den anden, er altså 0.4.
        
        \item Sandsynligheden for, at virksomheden ikke vinder nogen af kontrakterne, er sandsynligheden for komplementet til foreningsmængden af hændelserne $A$ og $B$. Ifølge komplementreglen er:
        
        \begin{align*}
            P((A \cup B)^c) &= 1 - P(A \cup B) \\
                                    &= 1 - 0.8 \\
                                    &= 0.2
        \end{align*}
        
        Sandsynligheden for, at virksomheden ikke vinder nogen af kontrakterne, er altså 0.2.
        
        \item Sandsynligheden for, at virksomheden vinder præcis én kontrakt, er sandsynligheden for differensen mellem foreningsmængden af hændelserne $A$ og $B$ og fællesmængden af hændelserne $A$ og $B$. Ifølge differensreglen er:
            
            \begin{align*}
                P((A \cup B) - (A \cap B)) &= P(A \cup B) - P(A \cap B) \\
                                            &= 0.8 - 0.2 \\
                                            &= 0.6
            \end{align*}
            Sandsynligheden for, at virksomheden vinder præcis én kontrakt, er altså 0.6.
            \end{enumerate}
\end{solution}

\begin{example}
    \begin{itemize}
    \item Der er 60 procents chance for, at det regner i dag.
    \item Der er 50 procents chance for, at det regner i morgen.
    \item Der er 30 procents chance for, at det ikke regner nogen af dagene.
    \end{itemize}
    
\begin{enumerate}[label=(\alph*)]
    \item Sandsynligheden for, at det regner i dag eller i morgen.
    \item Sandsynligheden for, at det regner i dag og i morgen.
    \item Sandsynligheden for, at det regner i dag, men ikke i morgen.
    \item Sandsynligheden for, at det regner enten i dag eller i morgen, men ikke begge dage.
\end{enumerate}
    
\begin{solution}
\begin{enumerate}[label=(\alph*)]
    \item     Lad \( A \) være hændelsen, at det regner i dag, og \( B \) hændelsen, at det regner i morgen.
        
    Vi har givet:
    \begin{align*}
        P(A \cup B) &= 1 - P\left((A \cup B)^c\right) \quad &\text{jf. komplementreglen} \\
                &= 1 - P\left(A^c \cap B^c\right) \quad &\text{jf. De Morgans lov} \\
                &= 1 - 0.3 \\
                &= 0.7
    \end{align*}
            
        
    Sandsynligheden for, at det regner i dag eller i morgen, er altså \( 0.7 \).

    \item Sandsynligheden for, at det regner i dag og i morgen: Det er $P(A \cap B)$. For at bestemme den bemærker vi, at
    \begin{align*}
    P(A \cap B) &= P(A) + P(B) - P(A \cup B) \\
                &= 0.6 + 0.5 - 0.7 \\
                &= 0.4
    \end{align*}

    \item Sandsynligheden for, at det regner i dag, men ikke i morgen: Det er $P\left(A \cap B^c\right)$.

    \begin{align*}
        P(A \cap B^c) &= P(A - B) \\
                      &= P(A) - P(A \cap B) \\
                      &= 0.6 - 0.4 \\
                      &= 0.2
        \end{align*}
        
    \item Sandsynligheden for, at det regner enten i dag eller i morgen, men ikke begge dage: Det er $P(A-B)+P(B-A)$. Vi har allerede fundet $P(A-B)=.2$. På tilsvarende måde kan vi bestemme $P(B-A)$:
    
    \begin{align*}
        P(B-A) & =P(B)-P(B \cap A) \\
        & =0.5-0.4 \\
        & =0.1
        \end{align*}
    Dermed er
    \begin{align*}
        P(A-B)+P(B-A) &= 0.2+0.1 \\
                      &= 0.3
    \end{align*}
\end{enumerate}
\end{solution}
\end{example}
